% Source printed pp. 416-420 / PDF pp. 424-428.
\section{射影 GIT}\label{src:8.2}
\begin{epigraph}
这不是数学, 这是神学.
\epigraphauthor{Paul Gordan}
\end{epigraph}
我们讨论 GIT 的射影情形: 若 $G$ 作用于 $G$-不变闭子概形 $U\subseteq\mathbb P(V)$, 则半稳定轨迹 (定义~\bookref{8.2.1}) 有射影良商 $U^{\mathrm{ss}}\sslash G$ (定理~\bookref{8.2.6}).
\subsection{半稳定性与稳定性}\label{src:8.2.1-subsection}
\begin{definition}\label{src:8.2.1}
设线性约化群 $G$ 作用于射影 $\Bbbk$-概形 $U$, $\mathcal O_U(1)$ 是带 $G$-作用的丰沛线丛. 定义\newterm{半稳定轨迹}{semistable locus}与\newterm{稳定轨迹}{stable locus}如下:\footnotemark
\[
U^{\mathrm{ss}}:=\left\{u\in U\ \middle|\ \begin{array}{l}\text{存在 }f\in\Gamma(U,\mathcal O_U(d))^G,\\d>0,\quad f(u)\ne0\end{array}\right\}.
\]
\[
U^{\mathrm s}:=\left\{u\in U\ \middle|\ \begin{array}{l}\text{存在 }f\in\Gamma(U,\mathcal O_U(d))^G,\ d>0,\ \text{使得}\\f(u)\ne0,\\\text{轨道 }Gu\subseteq U_f\text{ 闭, 且}\\\text{函数 }U\to\mathbb Z,\ x\mapsto\operatorname{dim} G_x\\\text{在 }u\text{ 的某个开邻域上为常数}\end{array}\right\}.
\]
\footnotetext{因为函数 $x\mapsto\dim G_x$ 上半连续, 当 $\dim G_u=0$ 时, 此条件自动成立.}
若 $u\in U^{\mathrm{ss}}$ (分别地, $u\in U^{\mathrm s}$), 则称点 $u\in U$ \textit{半稳定} (分别地, \textit{稳定}).\footnote{文献中有时将非半稳定点 $u\in U\setminus U^{\mathrm{ss}}$ 称为``不稳定''; 我们避免使用这一可能引起误解的术语.}
\end{definition}
\begin{remark}[\texorpdfstring{$G$}{G}-等变射影嵌入]\label{src:8.2.2}
若 $U$ 在 $\Bbbk$ 上固有, 则给定带 $G$-作用的丰沛线丛 $\mathcal O_U(1)$, 即甚丰沛的 $G$-线性化 (见 \S\bookref{B.1.31}), 等价于给定 $G$-等变嵌入 $U\hookrightarrow\mathbb P(V)$, 其中 $V$ 是有限维 $G$-表示. 要强调的是, 稳定和半稳定轨迹依赖于 $G$-等变嵌入 $U\hookrightarrow\mathbb P(V)$ 的选择: 不同嵌入给出不同的半稳定轨迹 $U^{\mathrm{ss}}$ 和不同的商 $U^{\mathrm{ss}}\sslash G$. 事实上, $U$ 上同一个甚丰沛线丛也可能具有不同的 $G$-作用和嵌入. 当 $G$ 半单时, $G$ 没有非平凡特征标 $G\to\mathbb G_m$, 因而 $U$ 上的线丛若存在 $G$-作用, 则此作用唯一. 另一方面, 当 $U$ 是正规射影簇时, Sumihiro 线性化定理 (\bookref{B.1.33}) 表明每个线丛 $L$ 都有带 $G$-作用的正张量幂 $L^{\otimes n}$. 例如, $\mathbb P^n$ 上的 $\mathcal O(1)$ 没有 $\mathrm{PGL}_{n+1}$-作用, 但 $\mathcal O(n+1)$ 有.
\end{remark}
\begin{remark}[非射影情形]\label{src:8.2.3}
若 $U$ 是任意 $\Bbbk$-概形 (或 $\Bbbk$-代数空间), $L$ 是带 $G$-作用的线丛, 则在定义~\bookref{8.2.1} 中对所用的每个不变截面 $f$ 额外要求补集 $U_f$ 仿射, 便可推广半稳定性和稳定性的定义. 若 $U$ 射影且 $\mathcal O_U(1)$ 丰沛, 则 $U_f$ 自动仿射.
\end{remark}
\begin{remark}[零锥]\label{src:8.2.4}
\newterm{零锥}{nullcone} $\widehat N\subseteq\mathbb A(V)$ 按定义是 $U\setminus U^{\mathrm{ss}}$ 上的仿射锥: 它是闭 $G$-不变轨迹, 由仿射锥中的点 $\widehat u\in\widehat U$ 组成, 这些点对 $\mathbb A(V)$ 上每个非常数 $G$-不变多项式均满足 $f(\widehat u)=0$.
\end{remark}
\subsection{GIT 第一基本定理}\label{src:8.2.2-subsection}
给定 $G$-等变闭子概形 $U\subseteq\mathbb P(V)$, 考虑分次环 $R=\bigoplus_{d\geq0}\Gamma(U,\mathcal O_U(d))$ 和态射
\begin{equation}\label{src:8.2.5}
\widetilde\pi:U^{\mathrm{ss}}\longrightarrow U^{\mathrm{ss}}\sslash G:=\Proj R^G.
\end{equation}
注意 $U^{\mathrm{ss}}$ 可以为空, 此时 $\Proj R^G$ 是空概形. 若 $U^{\mathrm{ss}}$ 非空, 它恰是有理映射 $\Proj R\dashrightarrow\Proj R^G$ 的定义区域.
\begin{theorem}[射影 GIT]\label{src:8.2.6}
设 $G$ 是代数闭域 $\Bbbk$ 上的线性约化群, $U\subseteq\mathbb P(V)$ 是 $G$-等变闭子概形, 其中 $V$ 是有限维 $G$-表示. 则有 Cartesian 图
\[
\begin{tikzcd}[ampersand replacement=\&]
U^{\mathrm s}\arrow[r,hook]\arrow[d]\arrow[dr,phantom,"\square"]\&U^{\mathrm{ss}}\arrow[r,hook]\arrow[d,"\widetilde\pi"]\&U\\
{U^{\mathrm s}/G}\arrow[r,hook]\&{U^{\mathrm{ss}}\sslash G}\&{}
\end{tikzcd}
\]
其中 $U^{\mathrm s}/G\subseteq U^{\mathrm{ss}}\sslash G$ 是开子概形, \bookeqref{8.2.5} 中的映射 $\widetilde\pi$ 是良商, 限制 $\widetilde\pi|_{U^{\mathrm s}}:U^{\mathrm s}\to U^{\mathrm s}/G$ 是几何商. 此外, $U^{\mathrm{ss}}\sslash G$ 射影, 并有丰沛线丛 $L$, 使对某个 $N$ 有 $\widetilde\pi^*L\cong\mathcal O_U(N)$.

若 $G$ 在 $U$ 上的作用一般具有有限稳定子, 则 $G$ 在 $U^{\mathrm s}$ 上的作用固有 (即作用映射 $G\times U^{\mathrm s}\to U^{\mathrm s}\times U^{\mathrm s}$ 固有), 换言之, $[U^{\mathrm s}/G]$ 分离.
\end{theorem}
\begin{proof}
因为 $U$ 射影, $R=\bigoplus_{d\geq0}\Gamma(U,\mathcal O_U(d))$ 在 $\Bbbk$ 上有限生成. 因而由推论~\bookref{7.5.13}(3), $R^G$ 也在 $\Bbbk$ 上有限生成, 所以 $U^{\mathrm{ss}}\sslash G=\Proj R^G$ 射影. 局部化与取不变量交换, 故对 $R^G$ 中每个正次数齐次元 $f$, 有 $(R^G)_{(f)}=(R_{(f)})^G$. 于是有 Cartesian 图
\[
\begin{tikzcd}[ampersand replacement=\&]
{U_f=\Spec R_{(f)}}\arrow[r,hook]\arrow[d]\arrow[dr,phantom,"\square"]\&U^{\mathrm{ss}}\arrow[r,hook]\arrow[d,"\widetilde\pi"]\&U\\
{U_f\sslash G=(U^{\mathrm{ss}}\sslash G)_f}\arrow[r,hook]\&{U^{\mathrm{ss}}\sslash G}.\&{}
\end{tikzcd}
\]
良商性质是 Zariski 局部的, 而 $(U^{\mathrm{ss}}\sslash G)_f$ 覆盖 $U^{\mathrm{ss}}\sslash G$, 故 $\widetilde\pi:U^{\mathrm{ss}}\to U^{\mathrm{ss}}\sslash G$ 是良商. 按构造, 存在整数 $N$, 使 $L:=\mathcal O_{U^{\mathrm{ss}}\sslash G}(N)$ 是丰沛线丛, 拉回为 $\mathcal O_U(N)|_{U^{\mathrm{ss}}}$.

为证明 $U^{\mathrm s}\to U^{\mathrm s}/G$ 是几何商, 只须证明 $U^{\mathrm s}$ 中每个 $G$-轨道均闭. 因为轨道退化时稳定子的维数增加, 实际上只须证明 $U^{\mathrm s}$ 中稳定子的维数局部常数. 按定义, 每个点 $u\in U^{\mathrm s}$ 都有开邻域 $V\subseteq U$, 使对所有 $v\in V$ 有 $\dim G_v=\dim G_u$. 因为 $\dim G=\dim G_v+\dim Gv$, 轨道维数在 $V$ 上为常数. 最后, 若 $U$ 有一个稠密开子集, 其上稳定子维数为 $0$, 则由稳定性的定义, 每个 $u\in U^{\mathrm s}$ 都有有限 (可能非约化的) 稳定子. 因为 $[U^{\mathrm s}/G]\to U^{\mathrm s}/G$ 也是良模空间, 且 $[U^{\mathrm s}/G]$ 具有拟有限对角, 命题~\bookref{7.5.35} 推出 $[U^{\mathrm s}/G]$ 分离.
\end{proof}
\subsection{初步例子}\label{src:8.2.3-subsection}
\begin{example}\label{src:8.2.7}
设 $\mathbb G_m$ 通过 $t\cdot[x:y:z]=[tx:t^{-1}y:z]$ 在 $\mathbb P^2$ 上作用. 半稳定轨迹是 $V(xy,z)=\{[0:1:0],[1:0:0]\}$ 的补集, 良商为 $(\mathbb P^2)^{\mathrm{ss}}\to\Proj\Bbbk[xy,z]=\mathbb P^1$. $xy=0$ 上的纤维是三个轨道之并, 它在 $(\mathbb P^2)^{\mathrm{ss}}$ 中的补集是稳定轨迹. 限制到 $z\ne0$ 上得到良商 $\mathbb A^2\to\mathbb A^1$, $(x,y)\mapsto xy$; $z=0$ 上的纤维则为无穷远直线除去 $[0:1:0]$ 和 $[1:0:0]$.
\end{example}
\begin{example}\label{src:8.2.8}
考虑 $\mathrm{SL}_2$ 在 $X=(\mathbb P^1)^4$ 上的对角作用和 $\mathrm{SL}_2$-等变 Segre 嵌入
\[
\begin{aligned}
(\mathbb P^1)^4&\longrightarrow\mathbb P^{15},\\
([x_1:y_1],\ldots,[x_4:y_4])&\longmapsto[x_1x_2x_3x_4:\cdots:y_1y_2y_3y_4].
\end{aligned}
\]
它对应于 $L:=\mathcal O(1)\boxtimes\cdots\boxtimes\mathcal O(1)$ 的 $\mathrm{SL}_2$-线性化. 不变量环 $\bigoplus_{d\geq0}\Gamma(X,L^{\otimes d})$ 由次数 $1$ 的广义交比
\begin{align*}
I_1&=(x_1y_2-x_2y_1)(x_3y_4-x_4y_3),\\
I_2&=(x_1y_3-x_3y_1)(x_2y_4-x_4y_2),\\
I_3&=(x_1y_4-x_4y_1)(x_2y_3-x_3y_2)
\end{align*}
生成, 它们满足线性关系 $I_1-I_2+I_3=0$. 不变量环为 $\Bbbk[I_1,I_2]$, 商为 $X^{\mathrm{ss}}\sslash\mathrm{SL}_2=\mathbb P^1$. 半稳定轨迹 $X^{\mathrm{ss}}$ 由至多两个点相同的元组组成, 稳定轨迹则由点两两不同的元组组成.
\end{example}
\bookfiguretarget{8.2.9}
\begin{figure}[H]
\centering
\begin{tikzcd}[ampersand replacement=\&,column sep=large,row sep=large]
{(\mathbb P^1)^{4,\mathrm s}}\arrow[d]\arrow[r,hook]\&{(\mathbb P^1)^{4,\mathrm{ss}}}\arrow[d]\\
{\mathbb P^1\setminus\{0,1,\infty\}}\arrow[r,hook]\&\mathbb P^1
\end{tikzcd}
\qquad
\begin{tikzcd}[ampersand replacement=\&,row sep=large]
{(p_1,p_2,p_3,p_4)}\arrow[d,mapsto]\\
\frac{(p_1-p_2)(p_3-p_4)}{(p_1-p_3)(p_2-p_4)}
\end{tikzcd}
\par\medskip
\begin{tikzcd}[ampersand replacement=\&,column sep=large,row sep=large]
{(0,0,\infty,\infty)}\arrow[d,mapsto]\&{(0,\infty,\infty,0)}\arrow[d,mapsto]\&{(0,\infty,0,\infty)}\arrow[d,mapsto]\\
0\arrow[r,dash]\&1\arrow[r,dash]\&\infty
\end{tikzcd}
\bookfigurehref{8.2.9}\caption{射影等价意义下的四个无序点}\label{src:8.2.9}
\end{figure}
点两两不同的有序元组 $(p_1,\ldots,p_4)$ 映到交比
\[
\frac{(p_1-p_2)(p_3-p_4)}{(p_1-p_3)(p_2-p_4)}.
\]
特别地, 两个稳定元组射影等价 (即属于同一 $\mathrm{SL}_2$-轨道), 当且仅当交比相同. 补集 $X^{\mathrm{ss}}\setminus X^{\mathrm s}$ 含有三个闭轨道: $(0,0,\infty,\infty)$, $(0,\infty,0,\infty)$ 和 $(0,\infty,\infty,0)$ 的 $\mathrm{SL}_2$-轨道. $(0,0,1,\infty)$ 或 $(1,\infty,0,0)$ 等元组在 $X^{\mathrm{ss}}$ 中的 $\mathrm{SL}_2$-轨道不闭, 其轨道闭包包含 $\mathrm{SL}_2\cdot(0,0,\infty,\infty)$. $\mathbb P^1$ 中 $n$ 个有序点的一般情形见例~\bookref{8.5.1}.
\subsection{其他 GIT 设置}\label{src:8.2.4-subsection}
\begin{exercise}[关于特征标的仿射 GIT]\label{src:8.2.10}
设 $U=\Spec A$ 是代数闭域 $\Bbbk$ 上的有限型概形, 线性约化群 $G$ 的作用由余作用 $\sigma:A\to\Gamma(G,\mathcal O_G)\otimes A$ 给定. 设 $\chi:G\to\mathbb G_m=\Spec\Bbbk[t]_t$ 是特征标, $\mathcal O_U(\chi)$ 是由 $\chi$ 确定的平凡线丛 $\mathcal O_U$ 上的 $G$-作用. 每个不变截面的补集 $U_f$ 都仿射, 故关于 $\mathcal O_U(\chi)$ 的半稳定性和稳定性定义 (定义~\bookref{8.2.1} 及注记~\bookref{8.2.3}) 化为:
\[
U^{\mathrm{ss}}:=\left\{u\in U\ \middle|\ \begin{array}{l}\text{存在 }f\in A,\ f(u)\ne0,\ \text{并对某个 }d>0\\\text{有 }\sigma(f)=\chi^*(t)^d\otimes f\end{array}\right\}.
\]
\[
U^{\mathrm s}:=\left\{u\in U\ \middle|\ \begin{array}{l}\text{存在 }f\in A,\ f(u)\ne0,\ \sigma(f)=\chi^*(t)^d\otimes f,\\d>0,\ \text{轨道 }Gu\subseteq U_f\text{ 闭, 且函数}\\x\mapsto\dim G_x\text{ 在 }u\text{ 的某个开邻域上为常数}\end{array}\right\}.
\]
定义 $U^{\mathrm{ss}}\sslash G:=\Proj\bigoplus_{d\geq0}A_d$, 其中 $A_d=\{f\in A\mid\sigma(f)=\chi^*(t)^d\otimes f\}$. 证明定理~\bookref{8.2.6} 的结论成立, 但 $U^{\mathrm{ss}}\sslash G$ 在 $A^G=A_0$ 上射影 (而非在 $\Bbbk$ 上射影).

例如, 对 $U=\mathbb A^n$ 上的伸缩 $\mathbb G_m$-作用和恒等特征标 $\chi=\mathrm{id}$, 有 $U^{\mathrm{ss}}=U^{\mathrm s}=\mathbb A^n\setminus0$, 商为 $\mathbb P^{n-1}$.
\end{exercise}
\begin{exercise}[仿射底上的射影 GIT]\label{src:8.2.11}
设 $\Bbbk$ 是代数闭域, $U$ 是仿射有限型 $\Bbbk$-概形 $\Spec B$ 上的射影概形. 设线性约化群 $G$ 相容地作用于 $U$ 和 $\Spec B$. 假定有 $G$-等变嵌入 $U\hookrightarrow\mathbb P(\mathcal E)$, 其中 $\mathcal E$ 是 $\Spec B$ 上带 $G$-作用的向量丛. 证明定理~\bookref{8.2.6} 的结论成立, 但 $U^{\mathrm{ss}}\sslash G$ 在 $B^G$ 上射影 (而非在 $\Bbbk$ 上射影).
\end{exercise}
\begin{remark}[辛约化]\label{src:8.2.12}
GIT 与辛几何有一种有趣的联系. 设 $G$ 是 $\mathbb C$ 上的约化群, 作用于光滑射影簇 $U\subseteq\mathbb P(V)$, 其中 $V$ 是 $n+1$ 维 $G$-表示. 设 $\omega$ 是 $U$ 上的辛形式, $K\subseteq G$ 是极大紧子群, $\mathfrak k$ 为其 Lie 代数. 存在\newterm{矩映射}{moment map}
\[
\mu:U\longrightarrow\mathfrak k^\vee,
\]
它关于 $\mathfrak k^\vee$ 上的余伴随作用 $K$-等变, 并满足 $d\mu(x)(\xi)\cdot a=\omega_x(\xi,v_x)$, 其中 $u\in U$, $\xi\in T_xU$, $a\in\mathfrak k$, $v_x$ 是 $K$ 在 $U$ 上的无穷小作用所给出的向量场. 则
\[
u\in U\text{ 半稳定}\quad\Longleftrightarrow\quad\overline{Gu}\cap\mu^{-1}(0)\ne\emptyset,
\]
且包含 $\mu^{-1}(0)\hookrightarrow U$ 诱导同胚 $\mu^{-1}(0)/K\to U^{\mathrm{ss}}\sslash G$. 见 \cite[\S8]{MFK94}.
\end{remark}
